\def\level{BTS}  % Classe à droite
\def\course{CIEL 2}
\def\chaptername{Lois continues} % Titre au centre
\def\docpart{Loi normale} % Partie du cours
\def\docname{C223}        % Identifiant à gauche

\pagestyle{cours_FP}
\makeheaderBTS{} % Affiche le bandeau

\begin{definition}[Loi normale]{}
	
 Une variable aléatoire X suit une loi normale d'espérance $\mu$ et d'écart-type $\sigma$, notée $\mathcal{N}(\mu;\sigma)$, lorsque la densité de probabilité $f$ est définie par:
\[f(x)=\dfrac{1}{\sigma\sqrt{2\pi}}e^{-\dfrac{1}{2}\left(\dfrac{x-\mu}{\sigma}\right)^2}\]

\begin{center}
\includegraphics[width = 0.9\linewidth]{loi_normale2.png}
\end{center}
\end{definition}

\begin{itemize}
	\item[$\bullet$] 	La courbe représentative de la densité de probabilité \og en cloche \fg{}  est appelée une courbe de Gauss. Elle admet un axe de symétrie d'équation $x=\mu$.

	\item[$\bullet$] L'adjectif « normale » s'explique par le fait que cette loi décrit et modélise la majorité des situations concrètes dans lesquelles le résultat dépend de nombreux paramètres aléatoires indépendants.

\end{itemize}

\begin{propriete}[Caractères de position et de dispersion]{}

	$X \sim \mathcal{N}(\mu;\sigma)$: Espérance $E(X)=\mu$ et écart-type $\sigma(X)=\sigma$.

\end{propriete}

\begin{minipage}{0.55\linewidth}

On remarque que:

\begin{itemize}
	\item[$\bullet$] 	$\mu$ est un caractère de position. Lorsque $\mu$ varie l'extremum de la courbe de déplace le long de l'axe des abscisses. 

	\item[$\bullet$] $\sigma$ est un caractère de dispersion. Lorsque $\sigma$ augmente, la répartition est plus étalée.

\end{itemize}
	


\end{minipage}
\begin{minipage}{0.45\linewidth}
\includegraphics[width = 0.9\linewidth]{musigma.png}
\end{minipage}


\bigskip

\begin{propriete}[Règle empirique]{}
Soit $X$ une variable aléatoire suivant la loi normale $\mathcal{N}(\mu;\sigma)$, alors:

\begin{minipage}{0.6\linewidth}
\[P(\mu-\sigma\le X\le\mu+\sigma)\approx 0,68\]

\[P(\mu-2\sigma\le X\le\mu+2\sigma)\approx 0,95\]

\[P(\mu-3\sigma\le X\le\mu+3\sigma)\approx 0,997\]
\end{minipage}
\begin{minipage}{0.4\linewidth}
\includegraphics[width=0.85\linewidth]{loi_normale3.png}
\end{minipage}
\end{propriete}


\pagebreak

\begin{easybox}[Approximation de la loi binomiale par la loi normale]{}

	
{\it Lorsque $n$ est grand et $p$ ni trop près de $0$, ni trop près de $1$, alors la loi binomiale est très proche d'une loi normale ayant même espérance et même écart-type.

En pratique, on convient de faire cette approximation pour $n \geq 50$, $p$ ni proche de $0$ ni de $1$ et $np(1 - p) > 10$.}


Dans les bonnes conditions, on peut approcher un loi binomiale $\mathscr{B}(n,p)$ par une loi normale $\mathscr{N}(\mu,\sigma)$ avec $\mu=np$ et $\sigma=\sqrt{np(1-p)}$


\begin{center}
	\includegraphics[width=0.99\linewidth]{approxim.png}
\end{center}

\end{easybox}


\begin{exercice_cours}[\quad]{Dans cette activité, chaque probabilité demandée sera arrondie à $10^{-3}$.}



Une enquête réalisée par un institut de sondage permet d'estimer que la probabilité qu'une lettre, choisie au hasard dans le courrier d'une entreprise, parvienne à son destinataire en France le lendemain est $0,7$.

Dans une des agences d'une grande entreprise, on admet que l'on expédie $500$ lettres par jour.

On note $X$ la variable aléatoire qui, à un jour tiré au hasard, associe le nombre de lettres qui parviendront à leur destinataire le lendemain.

On suppose que les acheminements de ces lettres se font en toute indépendance.

\begin{enumerate}
  \item 
  \begin{enumerate}
    \item Expliquer pourquoi $X$ suit une loi binomiale et préciser les paramètres de cette loi.
    
    \item Calculer l'espérance mathématique de $X$, puis la valeur arrondie à $10^{-2}$ de l'écart-type de $X$.
    
    \item Calculer la probabilité que $355$ lettres exactement, sur les $500$ expédiées un jour tiré au hasard, parviennent à leur destinataire le lendemain.
  \end{enumerate}

  \item Soit $P_0$ la probabilité qu'au moins $375$ des $500$ lettres expédiées un jour tiré au hasard parviennent à leur destinataire le lendemain. Calculer la probabilité $P_0$.

  \item Pour calculer une valeur approchée $P_0'$ de $P_0$, on pense approcher la loi de $X$ par la loi de $Y$, où $Y$ est une variable aléatoire suivant la loi normale $\mathscr{N}(350~;~10,25)$.
  
  \begin{enumerate}
    \item Justifier que les deux paramètres choisis sont pertinents.
    
    \item Calculer une valeur arrondie à $10^{-3}$ de $P_0'=P(Y\geq 375)$, avec $Y$ qui suit la loi normale $\mathscr{N}(350~;~10,25)$ puis comparer avec la valeur trouvée pour $P_0$.
  \end{enumerate}

  \item En utilisant l'approximation de la loi binomiale $X$ par la loi normale $Y$, estimer la probabilité que le nombre de lettres, sur les $500$ expédiées un jour choisi au hasard, parvenant à leur destinataire le lendemain, soit compris entre $340$ et $365$.
\end{enumerate}


\end{exercice_cours}
